% reference.tex
\documentclass[11pt, a4paper]{article}

\usepackage[T1]{fontenc}
\usepackage[utf8]{inputenc}
\usepackage{tgpagella}
\usepackage[spanish, es-noshorthands]{babel}
\usepackage{amsmath, amssymb, bm}
\usepackage[margin=2.5cm]{geometry}
\usepackage{fancyhdr}
\usepackage[most]{tcolorbox}

\pagestyle{fancy}
\fancyhf{}
\setlength{\headheight}{16pt}
\lhead{\textbf{UCB Sede Tarija} | Ing. Mecatrónica}
\rhead{Referencia: $\tau$ y Respuesta Exponencial (Ruta A)[cite: 4]}
\lfoot{Prof: M.Sc. Ing. Bernardo Quiroga Turdera[cite: 1, 3]}
\renewcommand{\headrulewidth}{0.4pt}
\definecolor{ucbblue}{RGB}{0, 51, 102}

\begin{document}

\begin{tcolorbox}[colback=ucbblue!5!white, colframe=ucbblue, title=\textbf{Artefacto 4: Referencia Exponencial (Condicional)}]
\end{tcolorbox}

\section*{1. Constante de Tiempo ($\tau$)}
Para un circuito RC equivalente:
$$ \tau = R_{eq} \cdot C $$[cite: 4]

\textbf{Interpretación Física:}
\begin{itemize}
    \item Un $\tau$ menor significa una respuesta \textbf{más rápida}[cite: 4].
    \item Un $\tau$ mayor significa una respuesta \textbf{más lenta}[cite: 4].
\end{itemize}

\section*{2. Forma General de Primer Orden}
$$ x(t) = x(\infty) + \left[ x(0^+) - x(\infty) \right] e^{-t/\tau} $$[cite: 4]

Antes de aplicar esta ecuación, se deben resolver las siguientes preguntas[cite: 4]:
\begin{enumerate}
    \item ¿Cuál es el estado inicial $x(0^+)$?[cite: 4]
    \item ¿Cuál es el estado final $x(\infty)$?[cite: 4]
    \item ¿Cuál es la constante de tiempo $\tau$?[cite: 4]
\end{enumerate}

\vspace{1cm}
\begin{tcolorbox}[colback=red!5!white, colframe=red!75!black, title=\textbf{ADVERTENCIA DE INGENIERÍA}]
    \centering
    Esta ecuación es una referencia de interpretación, \textbf{no} una hoja de memorización a ciegas[cite: 4].
\end{tcolorbox}

\end{document}
